\documentclass[11pt,a4paper]{article}
\usepackage[margin=25mm]{geometry}
\usepackage[T1]{fontenc}
\usepackage{amsmath,amssymb,amsthm,array,booktabs,tabularx}
\usepackage[hidelinks]{hyperref}
\setlength{\emergencystretch}{2em}
\allowdisplaybreaks
\title{Combinatorial lower bounds for torus grids and positive word grids}
\author{A complete proof of the requested lower bounds}
\date{2 November 2026}
\begin{document}
\maketitle

\begin{abstract}
We prove, with the conventions of the supplied problem brief, that
\[
\nu_{G_{a,k}}(t)=\frac{t}{2}\,\frac{(a-1)ka}{2}
\qquad(a\ge2,\ k\ge1,\ 0\le t\le2/a).
\]
We give a uniform, explicitly specified finite construction solving every
problem $P(a,k)$. Independently, we prove the factor bound directly for every
positive word with one-cycle permutation. This fulfills the alternative
expressly permitted in Target B and removes the need for the proposed
torus-word-to-torus-grid move equivalence. We also prove the subsequence
statement and the requested insertion assembly. Only the facts designated
[Lean] or [Written] in the brief are used as preexisting results. No published
grid-homology theorem or holomorphic result is an additional dependency.
\end{abstract}

\section{Conventions and a comparison principle}

All grid coordinates, rectangles, gradings and canonical classes have exactly
the conventions of Section 2 of the brief. In particular, marking coordinates
refer to squares, state coordinates to grid points, and the full differential
counts empty rectangles regardless of their $X$-markings. Integers modulo $n$
are compared using representatives in $\{0,\ldots,n-1\}$.

\textbf{Definition 1.1.}
For a nonzero chain $z$ and $0\le t\le2$, write
\[
L_t^G(z)=\max_{U^e x\in\operatorname{supp}z}
\left(\frac{t}{2}A_G(x)-e\right).
\]
Let $\mathcal Z(O)$ be the set of degree-zero cycles representing the canonical
class in the complex with $O$-markings $O$. This notation is legitimate:
the graded complex, its differential and its polynomial subcomplex depend
only on $O$. The canonical class, being the unique nonzero element of the
polynomial image, also depends only on $O$.

\textbf{Lemma 1.2 (finite minimum).}
For every grid of size at least two and $0\le t\le2$,
\[
\nu_G(t)=\min_{z\in\mathcal Z(O_G)}L_t^G(z)\in\mathbb R.
\]
\textit{Proof.}
There is at most one monomial on each state in degree zero, so this chain
group is finite-dimensional over the finite field $\mathbb F_2$. There are
therefore finitely many chains and finitely many canonical representatives.
This set is nonempty, and none of its elements is zero because the canonical
class is nonzero. By Fact 2.6.3, a chain is in the weighted filtration at level
$s$ exactly when each monomial in its support has level at most $s$.
Consequently the infimum in Definition 2.6.4 is the displayed finite
minimum. This also proves finiteness for every intermediate link grid used
below. $\square$

\textbf{Lemma 1.3 (comparison with a rescaled parameter).}
Let $G,Q$ be grids of the same size with the same $O$-markings. Suppose
$b>0$, $h\in\mathbb R$, and
\[
A_G(x)\ge h+bA_Q(x)\qquad\text{for every state }x.
\]
For every $t\ge0$ such that $t\le2$ and $bt\le2$,
\[
\nu_G(t)\ge\frac{th}{2}+\nu_Q(bt).
\]
\textit{Proof.}
For every integer $e$ and every state $x$,
\[
\lambda_t^G(U^e x)\ge\frac{th}{2}+\lambda_{bt}^Q(U^e x).
\]
The representative cycles in the two complexes are exactly the same set
$\mathcal Z(O)$. For each such cycle, taking the maximum over its support
preserves the inequality. Its $Q$-level is at least $\nu_Q(bt)$ by Lemma 1.2.
Taking the minimum of the $G$-levels proves the claim. This argument compares
different parameters directly; it does not invoke a same-parameter
filtered-map theorem. $\square$

\section{Auxiliary grids with zero threshold}

\textbf{Definition 2.1.}
For $n\ge2$ and any permutation $O$ of $\mathbb Z_n$, define
\[
Q(O):\qquad O_Q(c)=O(c),\qquad X_Q(c)=O(c-1\bmod n).
\]
This is a grid because $c\ne c-1$ and $O$ is injective. Denote by $e_c$ its
horizontal row arc, with $O$ in column $c$ and $X$ in column $c+1$.
All these arcs advance exactly one column cyclically.

\textbf{Lemma 2.2 (explicit reduction).}
Every $Q(O)$ is related by an explicitly determined finite sequence of grid
moves to
\[
B_2:\qquad O=(1,0),\quad X=(0,1).
\]
\textit{Proof.}
For each current size $n\ge3$, perform the following operations.

First cyclically translate rows to put $e_0$ in the bottom row. Move that row
upwards until the row immediately above it is $e_1$ or $e_{n-1}$, stopping at
the first such encounter. Every row passed has the form $e_j$ with
$2\le j\le n-2$. Its planar horizontal interval $[j,j+1]$ is disjoint from
$[0,1]$, including their endpoints. Each swap is therefore an allowed
ordinary row commutation. Since both neighboring edges occur somewhere
above the initial bottom row, this procedure stops.

If $e_1$ is above $e_0$, the adjacent rows, in columns $0,1$, have $O$ at the
southwest and northeast squares, $X$ at the southeast square, and a blank
northwest square. The other marking of the upper row is its $X$ in column
$2$. Delete column $1$ and the lower row by an $O$-destabilization, placing
the merged $O$ at the northwest square.

If $e_{n-1}$ is above $e_0$, cyclically translate columns by $+1$. Old columns
$n-1,0,1$ become columns $0,1,2$. The upper row now has $O$ in column $0$
and $X$ in column $1$, and the lower row has $O$ in column $1$ and $X$ in
column $2$. The southwest square of the block in columns $0,1$ is blank.
Delete column $1$ and the upper row by an $O$-destabilization, placing the
merged $O$ at that southwest square.

The blank square in either case is indeed unmarked: its row has its other
marking in column $2$, distinct from $0,1$ because $n\ge3$. Both
destabilizations merge the edges $0\longrightarrow1$ and
$1\longrightarrow2$. After deleting column $1$, their merged edge again
advances one column, and every other edge still advances one column
cyclically. Thus the new grid is again $Q(O')$ at size $n-1$.

For explicit relabelling, let $r_0,r_1$ be the rows of these two edges just
before destabilization, and let $\rho$ be the increasing relabelling after
deleting $r_0$. Then
\[
O'(0)=\rho(r_1),\qquad O'(c)=\rho(O(c+1))\quad(1\le c\le n-2).
\]
The remaining $X$-markings obey $X'(c)=O'(c-1)$, including column $0$,
whose $X$ is in row $\rho(O(n-1))=O'(n-2)$.

The size decreases once in each iteration. At size two, a cyclic row
translation if needed gives $B_2$. This specifies all moves and terminates
after $n-2$ destabilizations. When $n=3$, there are no intervening
nonneighboring edges, so the same procedure applies without commutations.
No commutation of rows sharing an endpoint is used. $\square$

\textbf{Lemma 2.3 (zero auxiliary threshold).}
For every $n\ge2$, every $O\in\mathfrak S_n$ and $0\le u\le2$,
\[
\nu_{Q(O)}(u)=0.
\]
\textit{Proof.}
The two states of $B_2$ have gradings
\[
\begin{array}{c|ccc}
x&M_O(x)&M_X(x)&A(x)\\
(0,1)&0&-1&0\\
(1,0)&-1&0&-1
\end{array}
\]
as follows by substitution into Definition 2.4.1 of the brief.
For each source state there are exactly two empty rectangles. From $(0,1)$,
both have $O$-count zero; from $(1,0)$, both have $O$-count one. Thus the two
terms cancel in each case over $\mathbb F_2$.
Only $(0,1)$ can carry an integer Laurent exponent in degree zero, and that
exponent is zero. It is the canonical anchor by Fact 2.7.3, and its level
is zero for all $u\in[0,2]$. Therefore $\nu_{B_2}(u)=0$.
Lemma 2.2 and the already-formalized native move invariance, Fact 2.8.2,
give the result. No topological identification of the auxiliary grids
with an unknot is required. $\square$

\textbf{Definition 2.4 (Alexander contribution).}
For any grid $G=(O,X)$ of size $n$, set
\[
C_G(c,r)=\sum_{j=0}^{c-1}
\left(\mathbf1_{\{X(j)<r\}}-\mathbf1_{\{O(j)<r\}}\right),
\qquad
K_G=\frac{J(O,O)-J(X,X)-(n-1)}{2}.
\]
The sum is empty when $c=0$.

\textbf{Lemma 2.5 (Alexander contribution formula).}
For every state $x$,
\[
A_G(x)=K_G+\sum_{c=0}^{n-1}C_G(c,x(c)).
\]
\textit{Proof.}
Subtracting the definitions of the two Maslov gradings gives
$A_G(x)=J_{\mathrm c}(x,X)-J_{\mathrm c}(x,O)+K_G$.
For a single state point $(c,r)$, twice the centered-$J$ difference is
\[
\sum_{j\ge c}\left(\mathbf1_{\{r\le X(j)\}}-\mathbf1_{\{r\le O(j)\}}\right)
+\sum_{j<c}\left(\mathbf1_{\{X(j)<r\}}-\mathbf1_{\{O(j)<r\}}\right).
\]
The first sum equals the second: both marking permutations have exactly
$n-r$ values at least $r$, and complementary inequalities are strict as
displayed. Hence the centered-$J$ difference is $C_G(c,r)$. Summing over
the state proves the formula, including boundary rows and columns.
$\square$

\textbf{Lemma 2.6 (auxiliary Alexander grading).}
For every $Q=Q(O)$ and every state $x$,
\[
A_Q(x)=-T_O(x),\qquad
T_O(x)=\#\{c:x(c)>O(c-1\bmod n)\}.
\]
\textit{Proof.}
Put $p=O(n-1)$ and $b_c=O(c-1)$. The sum defining $C_Q$ telescopes:
\[
C_Q(c,r)=\mathbf1_{\{p<r\}}-\mathbf1_{\{b_c<r\}}.
\]
For $c=0$ the right side is also zero. Since $x$ is a permutation,
\[
\sum_cC_Q(c,x(c))=(n-1-p)-T_O(x).
\]
The word of values of $X_Q$ is obtained from the word of values of $O$
by moving its last value $p$ to the front. Pairs not involving this value
are unchanged. The number of increasing pairs involving it changes from
$p$ to $n-1-p$. Thus
\[
J(X_Q,X_Q)-J(O,O)=n-1-2p,
\qquad K_Q=p-(n-1).
\]
Lemma 2.5 now gives $A_Q=-T_O$. $\square$

\section{The torus-grid lower bound: Target A}

\textbf{Definition 3.1.}
For $y\in\mathfrak S_N$, write
\[
F(y)=\sum_{c=0}^{N-1}|y(c)-c|,\qquad
D(y)=F(y)-\operatorname{inv}(y),\qquad
d(y)=\#\{c:y(c)<c\}.
\]
As in the brief, $\theta_s y$ adds $s$ to all values modulo $N$, and
$\bar x(c)=-x(c)\bmod N$.

\textbf{Lemma 3.2 (one value rotation).}
For all $N\ge2$ and $y\in\mathfrak S_N$,
\[
D(\theta_1y)-D(y)=N-1-2d(y).
\]
\textit{Proof.}
Let $p=y^{-1}(N-1)$. All value comparisons not involving column $p$
are unchanged by the rotation. The maximum at $p$ becomes the minimum:
its $N-1-p$ inversions with later columns disappear, and its $p$ pairs
with earlier columns become inversions. Hence
\[
\operatorname{inv}(\theta_1y)-\operatorname{inv}(y)=2p-(N-1).
\]
For $c\ne p$, the change in $|y(c)-c|$ is $-1$ when $y(c)<c$ and $+1$
otherwise. Column $p$ is not deficient, and its contribution changes by
$p-(N-1-p)=2p-(N-1)$. Therefore
\[
F(\theta_1y)-F(y)=(N-1-2d(y))+2p-(N-1)=2p-2d(y).
\]
Subtracting proves the claim. This includes $p=0$ and $p=N-1$; the
corresponding collections of inversion pairs are empty. $\square$

\textbf{Lemma 3.3 (rotation inequality).}
For every $N\ge2$, $y\in\mathfrak S_N$ and integer $1\le a<N$,
\[
D(\theta_a y)-D(y)\ge a(N-a)-2a\,d(y).
\]
\textit{Proof.}
For $0\le s<N$, exactly $s$ values wrap under addition of $s$ modulo $N$.
A nonwrapping value increases, so it cannot create a new deficiency.
Consequently $d(\theta_s y)\le d(y)+s$, including $s=0$.
Summing Lemma 3.2 for $\theta_s y$, $0\le s<a$, gives
\[
D(\theta_a y)-D(y)
=a(N-1)-2\sum_{s=0}^{a-1}d(\theta_s y)
\ge a(N-1)-2ad(y)-a(a-1).
\]
The last expression is $a(N-a)-2ad(y)$. All wrapping contributions
have been counted explicitly. $\square$

\textbf{Lemma 3.4 (the particular auxiliary grid).}
Define
\[
B_N:\qquad O(c)=N-1-c,\qquad X(c)=-c\bmod N.
\]
Then $B_N=Q(O)$, its threshold is zero on $[0,2]$, and
\[
A_{B_N}(x)=-d(\bar x)\qquad(x\in\mathfrak S_N).
\]
\textit{Proof.}
The equality $X(c)=O(c-1)$ is immediate, so Lemmas 2.3 and 2.6 apply.
To identify the count in Lemma 2.6, put $y=\bar x$ and let $p$ be the
column where $y(p)=x(p)=0$. For $c>0$ with $c\ne p$, the inequality
$x(c)>N-c$ is equivalent to $y(c)<c$. If $p>0$, column $p$ contributes
one to $d(y)$ and zero to $T_O(x)$, while column $0$ contributes zero
to $d(y)$ and one to $T_O(x)$. If $p=0$, neither exception contributes.
Thus $T_O(x)=d(y)$ in both cases.

For this particular family, the moves can be shortened to exactly
$N-2$ destabilizations. In columns $0,1$ and rows $N-2,N-1$, the
markings are $O$ northwest and southeast, $X$ northeast, and a blank
southwest square. For $N\ge3$, the other marking in column $0$ is
$X$ at row $0\ne N-2$, so the blank is genuine. Delete column $1$
and row $N-1$, placing $O$ at $(0,N-2)$. Old columns $c\ge2$ become
$c-1$, and the resulting permutations are exactly those of $B_{N-1}$.
$\square$

\textbf{Theorem 3.5 (M806, including the sharp endpoint).}
For every $a\ge2$, $k\ge1$, put
\[
q=ka+1,\qquad N=a+q,\qquad g=\frac{(a-1)(q-1)}{2}.
\]
Then
\[
\nu_{G_{a,k}}(t)=\frac{tg}{2}\qquad(0\le t\le2/a).
\]
In particular,
\[
\nu_{G_{a,k}}(2/a)=\frac{g}{a}=\frac{(a-1)k}{2}.
\]
\textit{Proof.}
For a state $x$, put $y=\bar x$. The written grading formulas
4.2.2 and 4.2.5 of the brief give
\[
A_{G_{a,k}}(x)=\frac{D(\theta_a y)-D(y)-(N-1)}{2}.
\]
By Lemmas 3.3 and 3.4,
\[
A_{G_{a,k}}(x)\ge
\frac{a(N-a)-(N-1)}{2}-a\,d(y)
=g+aA_{B_N}(x).
\]
Here the constant identity is
$a(N-a)-(N-1)=aq-a-q+1=(a-1)(q-1)=2g$.
The two grids have exactly the same $O$-markings. Apply Lemma 1.3
with $b=a$, $h=g$ and $Q=B_N$. Because $0\le at\le2$,
\[
\nu_{G_{a,k}}(t)\ge\frac{tg}{2}+\nu_{B_N}(at)=\frac{tg}{2}.
\]
The canonical anchor $x^O(c)=-c$ has $M_O=0$ and $A=g$ by Fact 2.7.3
and Lemma 4.2.5 of the brief. Its level is $tg/2$, proving the reverse
inequality. There is no restriction to $k=1$, and the proof applies
unchanged when $a=2$ or $t=0$. $\square$

\section{Uniform finite certificates and Target C}

This section makes the solution of $P(a,k)$ explicit as a finite formula.
The construction depends only on $N$, and the same set works for every
required parameter pair of that size.

\textbf{Definition 4.1 (allowed columns and rectangle matrix).}
Fix a required size $N=(k+1)a+1$ and set
\[
E_N=\{y\in\mathfrak S_N:D(y)\text{ is even and }D(y)\ge2d(y)\}.
\]
Order permutations lexicographically. Over $\mathbb F_2$, form the matrix
$R_N$ with columns $E_N$ and rows $\{*\}\cup\mathfrak S_N$, putting
the row $*$ first and ordering the other rows lexicographically. Define
\[
R_N(*,y)=\mathbf1_{\{y=\mathrm{id}\}},\qquad
R_N(v,y)=\#\operatorname{Rect}^{\circ}(y,v)\bmod2.
\]
Set $b_*=1$ and $b_v=0$ for every permutation row.

These entries are completely finite and use no marking restrictions.
More explicitly, for $u,v\in\mathbb Z_N$ let $\Delta(u,v)$ be the
representative of $v-u$ modulo $N$. For each ordered pair $L\ne R$,
include a contribution precisely when $v=y\circ(L\ R)$ and there is
no column $c$ satisfying both
\[
0<\Delta(L,c)<\Delta(L,R),\qquad
0<\Delta(y(L),y(c))<\Delta(y(L),y(R)).
\]
Sum these contributions modulo two. This formula includes rectangles
wrapping in either direction, as well as intervals with empty interior.

\textbf{Lemma 4.2 (consistency of the certificate equations).}
The system $R_N\eta=b$ is consistent.
\textit{Proof.}
By Lemma 3.4, $\nu_{B_N}(2)=0$. The dual certificate theorem, Fact 2.6.5,
therefore provides a functional on $C_0(B_N)$ that annihilates boundaries,
takes value one on the canonical anchor, and vanishes below level zero.
For $y=\bar x$ of even defect, Lemma 4.2.2 of the brief and Lemma 3.4 give
\[
M_O(x)=-D(y),\qquad
\lambda_2^{B_N}(U^{-D(y)/2}\bar y)=\frac{D(y)}{2}-d(y).
\]
Thus the functional's possible support consists exactly of columns in
$E_N$. Its value on the anchor gives the row $*$ equation.

For an odd-defect permutation $v$, the degree-one monomial on $\bar v$
is $U^{(-D(v)-1)/2}\bar v$. Reflection, Fact 2.8.6, puts its outgoing
empty rectangles to $\bar y$ in bijection with the rectangles from
$y$ to $v$. The Maslov rectangle formula, Fact 2.4.2, makes every
corresponding output exponent exactly $-D(y)/2$. Consequently the
equation that the functional annihilates this boundary is precisely
the row $v$ of $R_N\eta=b$.

For an even-defect $v$, all entries in that row vanish. Indeed,
$F(y)$ is even because the sum of positive displacements equals the
absolute sum of negative displacements. Hence $D(y)$ has the parity
of $\operatorname{inv}(y)$. A transposition changes that parity,
so an empty rectangle from an even-defect $y$ can only end at an
odd-defect state. This treats every row of the system. $\square$

\textbf{Definition 4.3 (an explicit determinant formula).}
Let $r$ be the largest order of a square minor of $R_N$ with determinant
one. Choose the lexicographically first pair of row and column index
sets $(I,J)$ of size $r$ with this property, and put $M=R_N[I,J]$.
There is at least one such minor since the anchor column has entry one
in row $*$. For $1\le j\le r$, let $M_j$ be obtained from $M$ by
replacing column $j$ with $b_I$. Define the vector $\eta_N$ by
\[
(\eta_N)_{J_j}=\det(M_j),\qquad
(\eta_N)_y=0\quad(y\notin J),
\qquad
Y_N=\{y\in E_N:(\eta_N)_y=1\}.
\]
All determinants and entries are in $\mathbb F_2$. This is a finite
formula indexed by $N$; Gaussian elimination is an equivalent
implementation.

\textbf{Lemma 4.4 (verification of the formula).}
The vector $\eta_N$ satisfies $R_N\eta_N=b$.
\textit{Proof.}
The columns indexed by $J$ are independent because their restriction
to $I$ is invertible. They span all columns: otherwise subtracting
from a column the linear combination matching it on $I$ would give
a nonzero residual in a row outside $I$, producing a nonzero minor
of order $r+1$. This contradicts the definition of $r$.
By Lemma 4.2, $b$ is in the column span, so it is a linear combination
of the columns in $J$. Restriction to $I$ determines those coefficients
uniquely. Cramer's rule, with $\det M=1$, gives precisely
$\det M_j$. Thus the full vector equation holds. $\square$

\textbf{Theorem 4.5 (all four conditions of $P(a,k)$).}
For every $a\ge2$, $k\ge1$, the set $Y_{(k+1)a+1}$ solves $P(a,k)$.
\textit{Proof.}
The anchor row of Lemma 4.4 gives $\mathrm{id}\in Y_N$. Membership
in $E_N$ gives even defect. For each target $v$, its row gives
\[
\sum_{y\in Y_N}\#\operatorname{Rect}^{\circ}(y,v)=0
\quad\text{in }\mathbb F_2.
\]
To spell out the rectangle cases: if $v$ is not obtained from $y$
by a transposition, its contribution is zero. If it is obtained by
swapping columns $L,R$, the only candidates are the two ordered
rectangles $(L,R)$ and $(R,L)$. Zero or two empty candidates
contribute zero; exactly one contributes one. Their emptiness is
tested by the modular-distance formula in Definition 4.1, with no
special exclusion for wrapping rectangles or marking contents.
The row equation states exactly that the sum of all these remaining
ones is even. Even-defect targets receive none, as proved in
Lemma 4.2. These are all possible rectangle cancellations.

Finally Lemma 3.3 gives, for $y\in Y_N$,
\[
(a-1)D(y)+D(\theta_a y)
\ge a(N-a)+a(D(y)-2d(y))\ge a(N-a)=aq.
\]
These are the four required conditions. $\square$

\textbf{Corollary 4.6 (functional and mirror cycle).}
An explicit dual functional is obtained by setting
\[
\varphi_{a,k}(U^{-D(\bar x)/2}x)=\mathbf1_{\{\bar x\in Y_N\}}
\]
on the degree-zero basis and extending linearly. It annihilates
boundaries by the preceding rectangle calculation and has value
one on the canonical anchor. Its support lies at or above $g/a$
at $t=2/a$. For any $0\le t\le2/a$, put $\sigma=at/2$. On this support,
\[
\lambda_t=(1-\sigma)\frac{D(\bar x)}{2}
+\sigma\lambda_{2/a}\ge\frac{tg}{2},
\]
using the written nonnegativity of $D$ from Lemma 4.2.4 of the brief.
Thus it certifies the entire interval, not only the endpoint.
The corresponding mirror cycle is explicitly
\[
z_{a,k}=\sum_{y\in Y_N}U^{D(y)/2}y\in C(H_{a,k}).
\]
It has degree $1-N$, coefficient one at $\mathrm{id}$ and nonnegative
exponents. For an empty rectangle $y\longrightarrow v$, the formula
$D(v)=D(y)-1+2|O_H\cap r|$ makes the output exponent
$(D(v)+1)/2$, independent of the source rectangle. The same parity
row therefore cancels its coefficient, proving $dz_{a,k}=0$.

\textbf{Corollary 4.7 (Target C).}
For every $k\ge1$, $P(2,k)$ is solved by $Y_{2k+3}$, and
\[
\nu_{G_{2,k}}(t)=\frac{tk}{2}\qquad(0\le t\le1).
\]
The argument uses the already-formalized native Fact 2.8.2, but none
of OSSbook Theorem 10.3.1, Proposition 4.6.15 or Chapter 5 as an
external theorem. In particular, it introduces none of the book
dependencies that Target C asks to remove.

\section{A direct proof for positive word grids}

We now establish Statement F directly. This also proves the
torus-word lower bound required by Target B's expressly permitted
alternative.

\textbf{Definition 5.1 (a preferred admissible word grid).}
Let $a\ge2$ and let $w=(s_0,\ldots,s_{\ell-1})$ be a positive word
in which every generator $0,\ldots,a-2$ occurs. Put $n=3\ell$ and
$\delta=(\ell-a+1)/2$. Construct its arcs as in Definition 2.11.3
of the brief, ordering rows by height and, within each height, by
increasing start column. Denote this grid by $W$.

\textbf{Lemma 5.2 (tie choices).}
Every admissible word grid of $w$ is related to $W$ by explicitly
specified row commutations.
\textit{Proof.}
At an integer strand height, write successive incoming and outgoing
columns as $I_j,O_j$. They satisfy
$I_j<O_j<I_{j+1}$, since the incoming and outgoing columns within
one occurrence differ by one and occurrences are ordered by letter
index. The arcs join $O_j$ to $I_{j+1}$ cyclically. The nonwrapping
planar intervals are pairwise disjoint. The unique wrapping
interval, drawn between its two endpoints in the planar cut,
strictly contains all the others. All endpoints are distinct.
Thus any two rows at this height are noninterleaving. If there is
only one occurrence, no row swap is needed.

At a half-integer height, each interval is $[3h,3h+2]$ for a
different letter index $h$, so these intervals are pairwise
disjoint. Again the corresponding rows are noninterleaving.
In each height block, perform adjacent swaps to sort the start
columns into increasing order, choosing the first adjacent
inversion at each step. Each is an allowed row commutation,
and the number of inversions decreases. This reaches $W$.
$\square$

\textbf{Definition 5.3 (active arcs, including wraparound).}
For an arc with start column $p$, end column $q$ and row rank $R$,
declare it active at the cut through grid points in column $c$ as
follows: if $p<q$, require $p<c\le q$; if $p>q$, require $c>p$
or $c\le q$. Equivalently, $c\in(p,q]_{\circ}$.
Let $F_c(r)$ count active arcs with rank less than $r$, and put
$U(r)=F_0(r)$. Thus $U(r)$ counts wrapping arcs with rank less
than $r$. This convention places the cut before the markings in
column $c$, exactly as required by the centered-$J$ grading.

\textbf{Lemma 5.4 (contribution from active arcs).}
For the word grid, and indeed for any grid with these directed
horizontal arcs,
\[
C_W(c,r)=U(r)-F_c(r).
\]
\textit{Proof.}
The contribution of one arc to the prefix sum in Definition 2.4 is
\[
\left(\mathbf1_{\{q<c\}}-\mathbf1_{\{p<c\}}\right)
\mathbf1_{\{R<r\}}.
\]
If $p<q$, its first factor is $-1$ for $p<c\le q$ and zero
otherwise. If $p>q$, it is $1$ for $q<c\le p$ and zero otherwise;
this is $1$ minus the active-arc indicator in that case.
Thus the factor is always the wrapping indicator minus the
active-arc indicator. Sum over all arcs. The inequalities also
cover $c=0$ and equality at either endpoint. $\square$

\textbf{Lemma 5.5 (active counts and newly started ranks).}
At every cut there are exactly $a$ active arcs, with distinct
heights. Put $b_c=O_W(c-1)$, and let $j_c$ be the number of active
arcs strictly below the active arc at row $b_c$. Then
\[
\sum_{c=0}^{n-1}j_c=\sum_{h=0}^{\ell-1}(3s_h+1).
\]
\textit{Proof.}
Immediately before a letter of index $s$, the active arcs are
one at each integer strand height $0,\ldots,a-1$. Across that
letter's three columns the successive replacements are:
the arc at $s$ is replaced by the middle arc at $s+1/2$;
the arc at $s+1$ is replaced by the outgoing arc at $s$;
the middle arc is replaced by the outgoing arc at $s+1$.
Thus the active count remains $a$ throughout and returns to
the integer heights before the next letter. At cut zero, the
active arcs are the unique wrapping connector at each integer
height, so this argument starts and closes cyclically.

The arc starting in column $c-1$ is active at cut $c$, and it
has row $b_c$. At cuts $3h+1,3h+2,3h+3$, its zero-based height
rank is respectively $s_h,s_h,s_h+1$. The final cut $n$ is cut
zero, so these triples account for all cuts exactly once.
Summing gives the formula. $\square$

\textbf{Lemma 5.6 (the two constant counts).}
For the preferred word grid,
\[
\sum_{r=0}^{n-1}U(r)=\sum_{h=0}^{\ell-1}(3s_h+2),
\qquad
J(O_W,O_W)-J(X_W,X_W)=2\ell-a.
\]
\textit{Proof of the first identity.}
At each integer height, the unique wrapping arc starts at the
last occurrence, so its row is last in its height class.
There are no wrapping arcs at half-integer heights. Therefore
a row of integer height $s$ has exactly $s$ wrapping rows below
it, and a row of height $s+1/2$ has exactly $s+1$ wrapping rows
below it. Each letter contributes its three starting arcs at
heights $s+1/2,s,s+1$, giving $(s+1)+s+(s+1)=3s+2$.

\textit{Proof of the second identity.}
Count increasing pairs of marking rows in column order.
For pairs from two different letter blocks with different
heights, the $O$ and $X$ contributions agree: each block has
the same height multiset $\{s,s+1/2,s+1\}$ for both markings.
Within a single block, the $O$ heights are $(s+1/2,s,s+1)$
and the $X$ heights are $(s,s+1,s+1/2)$; each has two
increasing pairs. Hence the total difference from unequal
heights is zero.

For equal half-integer heights, start and end columns of the
middle arcs have the same order, giving difference zero.
At an integer height with $m$ occurrences, the row ranks at
the $O$-columns occur in order $0,1,\ldots,m-1$, while the
row ranks at the $X$-columns occur in order
$m-1,0,\ldots,m-2$. The latter has $m-1$ fewer increasing
pairs. This also gives zero when $m=1$.
Every strand position occurs, and the total number of such
occurrences is $2\ell$. Summing $m-1$ over the $a$ integer
heights gives $2\ell-a$. These cases exhaust all pairs.
$\square$

\textbf{Theorem 5.7 (pointwise word-grid comparison).}
Let $Q=Q(O_W)$. For every state $x$,
\[
A_W(x)\ge\delta+aA_Q(x),\qquad
\delta=\frac{\ell-a+1}{2}.
\]
\textit{Proof.}
For each cut $c$, the active arc at row $b_c$ has exactly
$j_c$ active arcs below it. If $r\le b_c$, then
$F_c(r)\le j_c$. If $r>b_c$, then $F_c(r)\le a$ because
there are $a$ active arcs. Both cases give
\[
C_W(c,r)+a\mathbf1_{\{r>b_c\}}
\ge U(r)-j_c.
\]
Apply this with $r=x(c)$ and sum over columns. Since a state
uses every row exactly once, Lemmas 2.5, 2.6 and 5.6 give
\[
A_W(x)-aA_Q(x)
\ge K_W+\sum_{r=0}^{n-1}U(r)-\sum_{c=0}^{n-1}j_c.
\]
Here
\[
K_W=\frac{2\ell-a-(3\ell-1)}{2}
=\frac{-\ell-a+1}{2},
\]
and the difference of the two sums is $\ell$ by Lemmas 5.5
and 5.6. The resulting constant is
$(-\ell-a+1)/2+\ell=(\ell-a+1)/2=\delta$.
$\square$

\textbf{Theorem 5.8 (factor bound and Target B's alternative).}
For every $a\ge2$ and positive word $w$ containing every
generator, every admissible word grid satisfies
\[
\nu_{\operatorname{wordGrid}(w)}(t)
\ge\frac{t}{2}\,\frac{\ell(w)-a+1}{2}
\qquad(0\le t\le2/a).
\]
In particular, Statement F holds. Also, for every $a\ge2$
and $k\ge1$,
\[
\nu_{\operatorname{wordGrid}(c_a^{\,ka+1})}(t)
\ge\frac{tg}{2}\qquad(0\le t\le2/a).
\]
\textit{Proof.}
Apply Lemma 1.3 to Theorem 5.7 with $b=a$ and $h=\delta$.
Lemma 2.3 gives $\nu_Q(at)=0$, so the bound holds for $W$.
Lemma 5.2 and Fact 2.8.2 transfer it to every admissible
choice of word grid.

If the word permutation is an $a$-cycle, every generator
occurs: omission of generator $r$ would make every factor
preserve the nonempty proper subset $\{0,\ldots,r\}$, so
their product could not be an $a$-cycle. Moreover its sign
gives $\ell\equiv a-1\pmod2$, and full occurrence gives
$\ell\ge a-1$. Thus $\delta$ is the nonnegative integer
in Statement F, proving that statement.

For $c_a^q$ with $q=ka+1$, its length is $q(a-1)$ and
$\delta=(a-1)(q-1)/2=g$. Substitution proves the torus-word
bound. This is precisely the alternative to an explicit
move equivalence authorized in Target B: the lower bound
has been proved directly for the word grid itself. No
assertion of the separate move equivalence M807 is needed
or used. $\square$

\section{Subsequences and insertion assembly: Targets D and E}

\textbf{Proposition 6.1 (Target D).}
For $a\ge2$, a positive word $w=(s_0,\ldots,s_{\ell-1})$
and an integer $q\ge\ell$, the word $w$ is a subsequence
of $c_a^q$.
\textit{Proof.}
If $\ell=0$, use the empty embedding. Otherwise define
\[
f(i)=i(a-1)+s_i\qquad(0\le i<\ell).
\]
Then $0\le f(i)<(i+1)(a-1)\le q(a-1)$, and
\[
f(i+1)-f(i)=(a-1)+s_{i+1}-s_i
\ge(a-1)-(a-2)=1.
\]
Thus $f$ is strictly increasing and takes values in the
positions of $c_a^q$. Position $f(i)$ is letter $s_i$ in
the $i$-th Coxeter block. This is the required list
subsequence embedding. The case $q=0$ necessarily has
$\ell=0$ and was already covered. $\square$

\textbf{Lemma 6.2 (reversal of an $X$ switch).}
If $X'=X\circ\tau_c$ with $\tau_c=(c\ \ c+1\bmod n)$,
then $X=X'\circ\tau_c$.
\textit{Proof.}
Since $n\ge2$, the two columns are distinct and
$\tau_c^2=\mathrm{id}$. Compose the first equation with
$\tau_c$. This includes $c=n-1$, whose partner is column
zero, and requires no separate nonwrapping assumption.
$\square$

\textbf{Proposition 6.3 (Target E, with the required inputs explicit).}
Let $a\ge2$, let $w$ have one-cycle permutation, and suppose
$w$ embeds as a subsequence of $c_a^q$, where $q=ka+1$ and
$k\ge1$. Assume the written deletion statement M804.
For every $0\le t\le2/a$,
\[
\nu_{\operatorname{wordGrid}(w)}(t)
\ge\frac{t\delta(w)}{2}.
\]
This is Proposition 3.3.1 of the brief if its M806 and M807
assumptions are used; Theorem 5.8 supplies its needed torus-word
bound directly, so M807 can also be omitted.

\textit{Proof.}
Let $S$ be the image of the chosen subsequence embedding in
$\{0,\ldots,q(a-1)-1\}$. Write its complement in increasing
order as $j_1<\cdots<j_s$, where $s=q(a-1)-\ell(w)$.
Let $w_i$ be the restriction of $c_a^q$, in increasing
position order, to $S\cup\{j_1,\ldots,j_i\}$. Then
$w_0=w$, $w_s=c_a^q$, and each $w_{i+1}$ is obtained by
one insertion. Its insertion position is the number of
currently retained positions smaller than $j_{i+1}$.
Every $w_i$ contains $w$, so every generator still occurs.

Apply M804 to the insertion $w_i=u_iv_i$ and
$w_{i+1}=u_i\sigma_{r_i}v_i$. It supplies a grid $H_i$
obtained from the larger word grid by one $X$ switch,
with $H_i$ move-equivalent to the smaller word grid.
Lemma 6.2 reverses that switch. Facts 2.8.4 and 2.8.2
therefore give
\[
\nu_{\operatorname{wordGrid}(w_{i+1})}(t)
\le\nu_{\operatorname{wordGrid}(w_i)}(t)+\frac{t}{4}.
\]
The intermediate words may close to links, but all these
facts hold for their grids. Their sizes are $3\ell(w_i)\ge3$.

All thresholds here are finite real numbers by Lemma 1.2,
since $0\le t\le2/a\le1$ and the canonical classes are
nonzero. Hence one may sum and subtract these inequalities
in $\mathbb R$, and then view the result in the extended
reals. Summing, including the empty sum when $s=0$, gives
\[
\nu_{\operatorname{wordGrid}(c_a^q)}(t)
\le\nu_{\operatorname{wordGrid}(w)}(t)+\frac{st}{4}.
\]
The required torus-word lower bound is $tg/2$. It follows
either from the assumptions M806, M807 and Fact 2.8.2 in
the original proposition, or directly from Theorem 5.8.
Finally,
\[
s=q(a-1)-\ell(w)
=(a-1)(q-1)-2\delta(w)=2g-2\delta(w).
\]
Subtracting $st/4=t(g-\delta(w))/2$ gives the result.
For an arbitrary word, an admissible $k$ always exists;
for example $k=\max\{1,\ell(w)\}$ makes $ka+1\ge\ell(w)$,
so Proposition 6.1 applies. $\square$

\section{Boundary cases, dependencies and formalization interface}

\textbf{Boundary cases.}
The auxiliary grids are defined for every size at least two,
and every reduction stops at two. The first torus case is
$(a,k)=(2,1)$, with $N=5$, while the rotation argument works
for every $1\le a<N$. The case $k=1$ requires no separate
argument. At $a=2$ the interval is $[0,1]$ and the threshold
is $tk/2$. At $t=0$ both comparison parameters are zero;
at $t=2/a$ the auxiliary parameter is exactly two, where
the directly computed auxiliary threshold remains zero.
No assertion is made for $t>2/a$.

The empty word is included in the subsequence proposition.
It cannot have a one-cycle permutation on $a\ge2$ strands,
and it is not used to construct a word grid here. Every
word used for the word-grid proof contains all generators,
so each integer strand height has a nonempty occurrence
list and $n=3\ell\ge3$. Single-occurrence height classes
and intermediate links were covered explicitly. Reflection
is used only in constructing the dual certificate, with
the square and state conventions of Fact 2.8.6; the target
torus grid is always the reflected grid $G_{a,k}$.

\textbf{Dependency table.}
All citations in the following table refer to the supplied
brief, not to additional published results.

\begin{center}
\begin{tabularx}{\textwidth}{@{}p{0.26\textwidth}X >{\raggedright\arraybackslash}p{0.15\textwidth}@{}}
\toprule
New result & Permitted inputs from the brief & New published inputs\\
\midrule
Lemmas 1.2--1.3 & 2.5.2, 2.6.3, 2.7.1--2.7.3 & None\\
Lemmas 2.2--2.3 & Move definition 2.8.1; 2.7.3, 2.8.2 & None\\
Lemmas 2.5--2.6 & Grading definition 2.4.1 & None\\
Lemmas 3.2--3.4 & Permutation definitions; Section 2 above & None\\
Theorem 3.5 & 2.7.3; written 4.2.2, 4.2.5 & None\\
Lemmas 4.2--4.4 & 2.4.2, 2.6.5, 2.8.6; written 4.2.2 & None\\
Theorem 4.5 & Lemmas 3.3, 4.2--4.4 above & None\\
Corollaries 4.6--4.7 & Written 4.2.3--4.2.5; rectangle grading & None\\
Lemmas 5.2--5.6; Theorem 5.7 & 2.11.3; 2.4.1; move definition 2.8.1 & None\\
Theorem 5.8 & 2.8.2; Sections 1, 2, 5 above & None\\
Proposition 6.1 & Word and subsequence definitions & None\\
Proposition 6.3 & Written M804; 2.8.2, 2.8.4 & None\\
\bottomrule
\end{tabularx}
\end{center}

Native move invariance is a genuine dependency. It is used
exactly as the already-formalized Fact 2.8.2 permits. Neither
the stronger vanishing conjecture $V(a,k)$, nor an
identification with holomorphic knot Floer homology, nor
the separate move-equivalence statement M807 is assumed.
Target B is fulfilled through its allowed direct-word-grid
alternative, proved in Theorem 5.8.

\textbf{Minimal new combinatorial lemmas for formalization.}
The following finite statements suffice; all their proofs
are given above.

\begin{enumerate}
\item For $y\in\mathfrak S_N$, the exact identity
$D(\theta_1y)=D(y)+N-1-2\#\{c:y(c)<c\}$ and the bound
$\#\{c:(\theta_s y)(c)<c\}\le\#\{c:y(c)<c\}+s$ for
$0\le s<N$.
\item For every $O\in\mathfrak S_n$, the diagram
$(O,c\mapsto O(c-1))$ moves to $B_2$ by the explicit
size-decreasing algorithm of Lemma 2.2, and its Alexander
grading is $-\#\{c:x(c)>O(c-1)\}$.
\item For grids with equal $O$, equality of their graded
complexes and polynomial subcomplexes identifies their
canonical classes. A pointwise Alexander inequality
$A_G\ge h+bA_Q$ implies the rescaled threshold inequality
of Lemma 1.3.
\item The matrix with entries equal to empty-rectangle
cardinalities modulo two and an anchor row has the
explicit solution of Definition 4.3, after consistency
is obtained from the auxiliary dual certificate.
\item For the preferred admissible word grid, the
finite active-arc counts satisfy
$C(c,r)=U(r)-F_c(r)$, $\#\mathrm{Active}(c)=a$,
$\sum j_c=\sum_h(3s_h+1)$,
$\sum_rU(r)=\sum_h(3s_h+2)$, and
$J(O,O)-J(X,X)=2\ell-a$.
\item The strictly increasing list-position map
$i\mapsto i(a-1)+s_i$ proves the subsequence statement;
the adjacent-column transposition is its own inverse;
finite minima justify all threshold subtraction in
the insertion assembly.
\end{enumerate}

These statements prove Target A for every required pair,
give its full interval extension, prove Target C by
solving $P(2,k)$, fulfill Target B's stated alternative,
and prove Targets D and E. The direct positive-word proof
also establishes Statement F without using the insertion
route.

\end{document}
